% Generated draft: requires mathematical and source review.
\documentclass[11pt]{article}
\usepackage[margin=25mm]{geometry}
\usepackage{fontspec}
\setmainfont{Latin Modern Roman}
\usepackage{amsmath,amssymb,mathtools,unicode-math}
\setmathfont{Latin Modern Math}
\usepackage{enumitem,needspace,xurl,hyperref}
\setlength{\parindent}{0pt}
\setlength{\parskip}{5pt}
\setlength{\emergencystretch}{3em}
\newcommand{\N}{\mathbb{N}}
\newcommand{\Z}{\mathbb{Z}}
\newcommand{\Q}{\mathbb{Q}}
\newcommand{\R}{\mathbb{R}}
\newcommand{\C}{\mathbb{C}}
\newcommand{\F}{\mathbb{F}}
\newcommand{\A}{\mathbb{A}}
\newcommand{\PP}{\mathbb{P}}
\newcommand{\E}{\mathbb{E}}
\newcommand{\eps}{\varepsilon}
\newcommand{\dd}{\,\mathrm{d}}
\begin{document}
\section*{UnsolvedMath ID: 30002154}
\textbf{Source problem code:} catalogue entry
\subsection{Definitions and mathematical statement}
Let \(p\) be a prime. A \(p\)-divisible group \(D\) of height \(h\) over a scheme \(S\) is a sheaf of commutative groups for the fppf topology such that multiplication by \(p\) is surjective as a morphism of sheaves, \(D=\varinjlim_{n\geq 1}D[p^n]\), and each kernel \(D[p^n]=\ker([p^n]\colon D\to D)\) is represented by a finite locally free group scheme of rank \(p^{nh}\). Over a field \(k\) of characteristic \(p\), its dimension is \(d=\dim_k\operatorname{Lie}(D[p])\), and its codimension is \(c=h-d\).

\textbf{Classification objective (with unresolved scope).} For fixed dimension \(d\) and codimension \(c\), hence height \(h=d+c\), classify all \(p\)-divisible groups \(D\) in the intended base category. The requested classification also includes their specialization relations, homomorphisms at finite truncation levels, and useful invariants that are monotone under specialization. The source does not specify the base category or the equivalence relation for classification, so it does not determine a single precise classification problem.

For two groups \(D,E\) over a common base \(S\), one conventional meaning of homomorphisms at level \(n\geq 1\) is the abelian group
\[
\operatorname{Hom}_S\bigl(D[p^n],E[p^n]\bigr),
\]
whose elements are homomorphisms of \(S\)-group schemes. The source does not specify whether it requests these groups, the corresponding Hom functors on base schemes, or another finite-level construction. The phrase ``finite-level'' does not assert that these groups have finite cardinality.

A possible notion of specialization is supplied by a \(p\)-divisible group over a discrete valuation ring, with its generic and special fibers giving the two objects being compared, possibly after extending their residue fields. This is only an illustrative convention: the source does not specify the allowed families, fields, or base changes. Likewise, specialization-monotonicity requires a specified specialization relation and an ordered target for each invariant, together with a direction of inequality. These choices remain unresolved. The definitions above clarify standard terminology without selecting missing hypotheses for the research objective.
\subsection{Short English statement}
For fixed dimension and codimension, the supplied problem asks for a classification of \(p\)-divisible groups, their possible specializations, their homomorphisms at finite levels, and useful invariants monotone under specialization. The base category, classification equivalence, and precise meanings of the additional requests are unspecified. The supplied record labels the problem open; this statement-only draft requires mathematical and source review.
\subsection{Sources}
\begin{itemize}
\item Supplied local UnsolvedMath registry record: permanent ID \(30002154\), title ``Classification and Specialization of p-Divisible Groups,'' source problem code \texttt{catalogue entry}, in the Oberwolfach Reports Open Problems collection. No individual proposer is named; the supplied proposed year is 2012.
\item Supplied source citation: \emph{Arithmetic Geometry (2013)}, Oberwolfach Reports, pp.\ 2335--2388, DOI \texttt{10.4171/owr/2012/38}. \url{https://doi.org/10.4171/owr/2012/38}. The citation's year 2013 differs from the year 2012 in the DOI and supplied proposed year; this discrepancy is unresolved here.
\item Supplied literature source labeled ``Bellovin, cohomology of phi-Gamma modules'': \url{https://arxiv.org/abs/2102.04820}. Its relevance to the stated classification of \(p\)-divisible groups is not established by the supplied record.
\item Supplied literature source labeled ``Oberwolfach report'': \url{https://doi.org/10.4171/owr/2012/38}.
\item Literature assessment dated exactly \texttt{2026-08-21}: the supplied curation labels the problem open and states that classifying all \(D\) is an explicit research objective, while existing work classifies rank-one modules and develops families rather than all objects. This is supplied dated curation, not an independently verified literature assessment. The record itself requests independent verification of time-sensitive status and cited claims.
\item The supplied statement status is \texttt{corrected\_verified}. Its verification note reports that the official report was checked, that it confirms the complete problem statement, and that prime, group, and Hom-functor notation were restored without changing the classification problem. These are supplied verification claims; no source visit or independent verification is asserted in this draft.
\end{itemize}
\end{document}
